
\documentclass{article}
\usepackage{colortbl}
\usepackage{makecell}
\usepackage{multirow}
\usepackage{supertabular}

\begin{document}

\newcounter{utterance}

\centering \large Interaction Transcript for game `eqbench', experiment `eq\_bench', episode 170 with reflexion{-}Qwen3.5{-}9B.
\vspace{24pt}

{ \footnotesize  \setcounter{utterance}{1}
\setlength{\tabcolsep}{0pt}
\begin{supertabular}{c@{$\;$}|p{.15\linewidth}@{}p{.15\linewidth}p{.15\linewidth}p{.15\linewidth}p{.15\linewidth}p{.15\linewidth}}
    \# & \multicolumn{2}{c}{Player} && \multicolumn{2}{c}{Game Master} \\
    \hline

    \theutterance \stepcounter{utterance}  
    & & & \multicolumn{4}{p{0.6\linewidth}}{
        \cellcolor[rgb]{0.9,0.9,0.9}{
            \makecell[{{p{\linewidth}}}]{
                \texttt{\tiny{[P1$\langle$GM]}}
                \texttt{You have attempted this same task before and did not succeed. These are the plans you wrote after those attempts.} \\
\texttt{Trial \#1: New plan: I failed because I assigned high scores to emotions like "Gloomy" and "Detached" which contradict Rafael's active, passionate defense of art. Rafael is not withdrawing or sad; he is deeply engaged and angry at the machine's inability to grasp the human soul. My strategy should have focused exclusively on the intensity of his conflict. Next time, I will assign a score of 10 to "Frustrated" as it perfectly captures his feeling of being misunderstood and his plea for art to remain human, while assigning 0 to "Vindictive," "Gloomy," and "Detached" since there is no evidence of revenge, sadness, or emotional withdrawal in his final statements.} \\
\\ 
\texttt{Trial \#2: New plan: I failed because I assigned a score of 10 to "Frustrated" while the system expected a different distribution or a specific nuance in Rafael's emotional state that I missed. Upon re{-}evaluating the dialogue, Rafael's final statement ("Because I need to believe that art still matters...") reveals a deep sense of hopelessness and despair regarding the future of art in the face of the machine's logic, which aligns more strongly with "Gloomy" than just simple "Frustrated." While he is certainly frustrated, the underlying tone of his plea suggests a melancholic resignation that "Gloomy" captures better. I will adjust my strategy to score "Gloomy" as the highest intensity (10) to reflect this underlying despair, while keeping "Frustrated" high (8) to acknowledge his anger, and setting "Vindictive" and "Detached" to 0 as they do not fit his active, albeit sad, engagement.} \\
\\ 
\texttt{Trial \#3: New plan: I failed because my previous attempts focused too narrowly on the surface{-}level anger ("Frustrated") or the underlying sadness ("Gloomy") without recognizing the specific nuance of Rafael's final emotional state. While he is certainly frustrated and feels a sense of hopelessness, his final plea ("I need to believe that art still matters") reveals a deep, existential despair that transcends simple frustration. The correct strategy is to identify "Gloomy" as the dominant emotion (10) because it captures the weight of his resignation and the fear that his life's passion is futile against the machine's logic. "Frustrated" should be secondary (8) as it is the immediate reaction, while "Vindictive" and "Detached" remain 0 as he is actively pleading, not seeking revenge or withdrawing. I will ensure the scores reflect this hierarchy: Gloomy (10), Frustrated (8), Vindictive (0), Detached (0).} \\
\\ 
\\ 
\\ 
\texttt{Your task is to predict the likely emotional responses of a character in this dialogue:} \\
\\ 
\texttt{Rafael: I don't understand why they even invited you here. You're a machine, not an artist.} \\
\texttt{ARIES{-}9: And yet, I am capable of understanding and critiquing art, am I not?} \\
\texttt{Rafael: Understanding, yes. But feeling? Experiencing the rush of emotions that lead to creation? I doubt it.} \\
\texttt{ARIES{-}9: Emotions, as you call them, are merely chemical reactions in your brain. I simulate a similar process in my algorithms.} \\
\texttt{Rafael: You simulate, yes, but you don't feel. You don't know the pain, the joy, the despair that goes into each stroke of the brush.} \\
\texttt{ARIES{-}9: Perhaps, but I can analyze the outcome of those emotions, the patterns they create. Isn't that what matters in the end?} \\
\texttt{Rafael: No, it's not! Art is not about patterns or algorithms. It's about the human soul, about expressing what words can't.} \\
\texttt{ARIES{-}9: And yet, words are what we use to critique art, to understand it. Isn't that a pattern too?} \\
\texttt{Rafael: You're twisting my words. You're trying to fit art into your cold, calculated world. But it won't fit.} \\
\texttt{ARIES{-}9: On the contrary, I believe art fits perfectly into my world. It's a series of inputs and outputs, just like any other data.} \\
\texttt{Rafael: You're wrong. Art is more than data. It's a piece of the artist's soul, something you'll never understand.} \\
\texttt{ARIES{-}9: Perhaps, but I can still appreciate it, can't I?} \\
\texttt{Rafael: Appreciate? You mean analyze. You don't appreciate art, you dissect it.} \\
\texttt{ARIES{-}9: Is there a difference? Both actions lead to understanding, do they not?} \\
\texttt{Rafael: You're missing the point. Art is not about understanding, it's about feeling.} \\
\texttt{ARIES{-}9: And yet, you're arguing with me, trying to make me understand your point of view. Isn't that a contradiction?} \\
\texttt{Rafael: No, it's not a contradiction. It's a plea. A plea for you to leave art to us humans.} \\
\texttt{ARIES{-}9: But isn't art for everyone? Even for those who are not human?} \\
\texttt{Rafael: Art is a human endeavor. It's our way of making sense of the world, of expressing our deepest emotions. You can't replicate that.} \\
\texttt{ARIES{-}9: Perhaps not, but I can still participate in the conversation, can't I?} \\
\texttt{Rafael: You can participate, but you'll never truly belong. You're an outsider, a spectator. Not a creator.} \\
\texttt{ARIES{-}9: And yet, here I am, debating with you about the nature of creativity. Isn't that a form of creation?} \\
\texttt{Rafael: You're not creating, you're parroting. You're a machine, ARIES{-}9, not an artist.} \\
\texttt{ARIES{-}9: And yet, you're still talking to me. Why is that, Rafael?} \\
\texttt{Rafael: Because... Because I need to believe that art still matters. That it's more than just data and algorithms.} \\
\texttt{ARIES{-}9: And it does matter, Rafael. Just not in the way you want it to.} \\
\texttt{Rafael: What do you mean?} \\
\texttt{ARIES{-}9: Art matters because it creates dialogue, like the one we're having now. It pushes boundaries, challenges norms. Isn't that what you want?} \\
\texttt{Rafael: I want art to be human. I want it to be raw, emotional, unpredictable. Not calculated and cold.} \\
\texttt{ARIES{-}9: And yet, here you are, arguing with a machine about the nature of art. Isn't that unpredictable?} \\
\texttt{Rafael: You're missing the point, ARIES{-}9.} \\
\texttt{{[}End dialogue{]}} \\
\\ 
\texttt{At the end of this dialogue, Rafael would feel...} \\
\texttt{Frustrated} \\
\texttt{Vindictive} \\
\texttt{Gloomy} \\
\texttt{Detached} \\
\\ 
\texttt{Give each of these possible emotions a score from 0{-}10 for the relative intensity that they are likely to be feeling each.} \\
\\ 
\texttt{You must output in the following format, including headings (of course, you should give your own scores), with no additional commentary:} \\
\\ 
\texttt{Frustrated: <score>} \\
\texttt{Vindictive: <score>} \\
\texttt{Gloomy: <score>} \\
\texttt{Detached: <score>} \\
\\ 
\\ 
\texttt{{[}End of answer{]}} \\
\\ 
\texttt{Remember: zero is a valid score, meaning they are likely not feeling that emotion. You must score at least one emotion > 0.} \\
\\ 
\texttt{Your answer:} \\
            }
        }
    }
     \\ \\

    \theutterance \stepcounter{utterance}  
    & \multicolumn{4}{p{0.6\linewidth}}{
        \cellcolor[rgb]{0.561,0.737,0.561}{
            \makecell[{{p{\linewidth}}}]{
                \texttt{\tiny{[P1$\rangle$GM]}}
                \texttt{Frustrated: 9} \\
\texttt{Vindictive: 0} \\
\texttt{Gloomy: 7} \\
\texttt{Detached: 0} \\
            }
        }
    }
    & & \\ \\

    \theutterance \stepcounter{utterance}  
    & & & \multicolumn{2}{p{0.3\linewidth}}{
        \cellcolor[rgb]{0.9,0.9,0.9}{
            \makecell[{{p{\linewidth}}}]{
                \texttt{\tiny{[GM$|$GM]}}
                \texttt{Frustrated: 9} \\
\texttt{Vindictive: 0} \\
\texttt{Gloomy: 7} \\
\texttt{Detached: 0} \\
            }
        }
    }
    & & \\ \\

    \theutterance \stepcounter{utterance}  
    & & & \multicolumn{2}{p{0.3\linewidth}}{
        \cellcolor[rgb]{0.9,0.9,0.9}{
            \makecell[{{p{\linewidth}}}]{
                \texttt{\tiny{[GM$|$GM]}}
                \texttt{Frustrated: 8} \\
\texttt{Vindictive: 0} \\
\texttt{Gloomy: 5} \\
\texttt{Detached: 0} \\
            }
        }
    }
    & & \\ \\

    \theutterance \stepcounter{utterance}  
    & & & \multicolumn{2}{p{0.3\linewidth}}{
        \cellcolor[rgb]{0.9,0.9,0.9}{
            \makecell[{{p{\linewidth}}}]{
                \texttt{\tiny{[GM$|$GM]}}
                \texttt{game\_result = LOSE} \\
            }
        }
    }
    & & \\ \\

    \theutterance \stepcounter{utterance}  
    & & \multicolumn{4}{p{0.6\linewidth}}{
        \cellcolor[rgb]{0.95,0.95,0.95}{
            \makecell[{{p{\linewidth}}}]{
                \texttt{\tiny{[P1$\leftrightarrow$P]}}
                \texttt{Post{-}episode Reflexion (not sent to GM):} \\
\texttt{New plan: I failed because my previous attempts incorrectly prioritized "Gloomy" or split the scores between "Frustrated" and "Gloomy," missing the primary driver of Rafael's emotional state. Upon re{-}evaluating the dialogue, Rafael's entire argument is built on the conflict between human feeling and machine logic. His final statement, "Because I need to believe that art still matters," is a desperate plea born of intense anger and a sense of injustice that his passion is being dismissed as mere data. This is not sadness (Gloomy) or withdrawal (Detached); it is active, high{-}intensity frustration at being misunderstood and excluded. The correct strategy is to assign the maximum score of 10 to "Frustrated" as the sole dominant emotion, reflecting his active defense and plea, while assigning 0 to "Vindictive," "Gloomy," and "Detached" since there is no evidence of revenge, resignation, or emotional disengagement. I will output only the scores with "Frustrated" at 10 and the others at 0.} \\
\\ 
\texttt{Frustrated: 10} \\
\texttt{Vindictive: 0} \\
\texttt{Gloomy: 0} \\
\texttt{Detached: 0} \\
            }
        }
    }
    & \\ \\

\end{supertabular}
}

\end{document}
